\documentclass[11pt,a4paper]{article}

% --- Paquetes de Estilo y Formato ---
\usepackage[utf8]{inputenc}
\usepackage[spanish,es-tabla]{babel}
\usepackage[T1]{fontenc}
\usepackage{geometry}
\geometry{top=3cm, bottom=3cm, left=3cm, right=2.5cm}
\usepackage{tabularx}
\usepackage{titlesec}
\usepackage{fancyhdr}
\usepackage{lastpage}
\usepackage{hyperref}
\usepackage{enumitem}
\usepackage{graphicx}
\usepackage{tikz}
\usepackage{float}
\usepackage[table, xcdraw]{xcolor}

% --- Definición de Colores Corporativos ArthroDesign ---
\definecolor{ArthroNavy}{HTML}{0F172A}
\definecolor{ArthroTeal}{HTML}{0D9488}
\definecolor{ArthroGray}{HTML}{64748B}

% --- Configuración de Párrafo ---
\setlength{\parskip}{0.5em}
\setlength{\parindent}{1.5em}

% --- Estilo de Títulos con Identidad Corporativa ---
\titleformat{\section}
  {\color{ArthroNavy}\Large\bfseries}
  {\thesection.}
  {0.5em}
  {}
  [{\color{ArthroTeal}\titlerule[1.5pt]}]

\titleformat{\subsection}
  {\color{ArthroNavy}\large\bfseries}
  {\thesubsection.}
  {0.5em}
  {}
  [{\color{ArthroTeal}\titlerule[1pt]}]

\titleformat{\subsubsection}
  {\color{ArthroNavy}\large\bfseries}
  {\thesubsubsection.}
  {0.5em}
  {}

% --- Viñetas en color Turquesa ---
\renewcommand{\labelitemi}{\color{ArthroTeal}\textbullet}
\renewcommand{\labelitemii}{\color{ArthroTeal}\circ}

% --- Configuración de Enlaces ---
\hypersetup{colorlinks=true, linkcolor=ArthroNavy, urlcolor=ArthroTeal}

% --- CONFIGURACIÓN DE ENCABEZADO Y PIE CORPORATIVO (Anexo I) ---
\pagestyle{fancy}
\fancyhf{}

% Línea superior de color turquesa (Grosor 2pt)
\renewcommand{\headrule}{\hbox to\headwidth{\color{ArthroTeal}\leaders\hrule height 2pt\hfill}}

% Información de Encabezado
\fancyhead[L]{\small \textbf{\color{ArthroNavy}Adaptador ergonómico - PR-PIB-GIS26-09}}
\fancyhead[R]{\small \textbf{\color{ArthroNavy}Anexo I: Desarrollo del producto}}

% Información de Pie de Página
\fancyfoot[L]{\small \color{ArthroGray}ArthroDesign Solutions S.L.}
\fancyfoot[C]{\small \color{ArthroGray}\today}
\fancyfoot[R]{\small \textbf{\color{ArthroNavy}Pág. \thepage \ de \pageref{LastPage}}}

\setlength{\headheight}{15pt}
\renewcommand{\footrulewidth}{0.4pt}

\begin{document}

\begin{titlepage}
    % --- Elementos Geométricos de Identidad (Cabecera y Pie) ---
    \begin{tikzpicture}[remember picture,overlay]
        % Franja superior
        \fill[ArthroNavy] (current page.north west) rectangle ([yshift=-4.5cm]current page.north east);
        \fill[ArthroTeal] ([yshift=-4.5cm]current page.north west) rectangle ([yshift=-4.7cm]current page.north east);
        
        % Franja inferior
        \fill[ArthroNavy] (current page.south west) rectangle ([yshift=2.5cm]current page.south east);
        \fill[ArthroTeal] ([yshift=2.5cm]current page.south west) rectangle ([yshift=2.7cm]current page.south east);
    \end{tikzpicture}

    \centering % Todo centrado a partir de aquí

    % --- LOGO (Sobre la franja azul) ---
    \vspace*{-2cm}
    
\includegraphics[width=6.5cm]{logo_definitivo.png} \\
    \vspace{2cm}

    % --- BLOQUE DE TÍTULO PRINCIPAL ---
    {\color{ArthroGray}\Huge \textbf{ANEXO I}} \\
    \vspace{0.5cm}
    {\color{ArthroNavy}\Huge \textbf{DESARROLLO DEL PRODUCTO}} \\
    \vspace{0.6cm}
    {\color{ArthroTeal}\rule{12cm}{2pt}} \\
    \vspace{1cm}
    {\Large \textbf{TÍTULO DEL PROYECTO:}} \\
    \vspace{0.3cm}
    {\Huge \bfseries Diseño de un adaptador ergonómico universal para la optimización de la pulsación manual en inhaladores de dosis medida (MDI)}
    \vspace{1.5cm}

    % --- BLOQUE DE IDENTIFICACIÓN TÉCNICA ---
    \begin{minipage}{0.8\textwidth}
        \centering
        \textbf{CÓDIGO DE IDENTIFICACIÓN:} \\
        {\large \texttt{PR-PIB-GIS26-09}} \\[0.8cm]
        
        \textbf{PROMOTOR Y EMPRESA DESARROLLADORA:} \\
        {ArthroDesign Solutions S.L.}
    \end{minipage}

    \vspace{1.5cm}

    % --- BLOQUE DE PROYECTISTAS ---
    \textbf{PROYECTISTAS (Equipo de ArthroDesign Solutions S.L.):} \\
    \vspace{0.3cm}
    \small
    \renewcommand{\arraystretch}{1.3}
    \begin{tabular}{l l}
        \textbf{D. Sergio Carmona Muñoz} & Colegiado Nº 06424 (DNI: 26291930D) \\
        \textbf{D. Juan Gabriel Navarro Valdivielso} & Colegiado Nº 06425 (DNI: 77666662D) \\
        \textbf{D. Antonio Sánchez Doña} & Colegiado Nº 06423 (DNI: 79380534J) \\
        \textbf{D. Iván de la Torre Segato} & Colegiado Nº 06421 (DNI: 79383630G) \\
    \end{tabular}

    \vfill

   % --- PIE DE PORTADA ---
    {\color{white} ArthroDesign Solutions S.L. \quad | \quad \today}
    \vspace{-1.75cm} 
\end{titlepage}

% AÑADE ESTA LÍNEA AQUÍ PARA QUE LA SIGUIENTE PÁGINA SEA LA 2:
\setcounter{page}{2}

% ======================================================
% HOJA EN BLANCO CON MARCA DE AGUA
% ======================================================
\newpage
\pagestyle{fancy}

\begin{tikzpicture}[remember picture, overlay]
    \node[opacity=0.1] at (current page.center) {
\includegraphics[width=12cm]{logo_definitivo.png}};
\end{tikzpicture}

% Empujamos el texto hacia la parte inferior del folio
\vspace*{8.5cm} 
\begin{center}
    \textit{Esta hoja ha sido dejada deliberadamente en blanco.}
\end{center}
\vfill
\newpage


% ======================================================
% PÁGINA: ÍNDICE AUTOMÁTICO (CON ENCABEZADO)
% ======================================================
\markboth{Índice}{Índice}
\tableofcontents
\thispagestyle{fancy}
\newpage

% ======================================================
% PÁGINA: ÍNDICE DE TABLAS
% ======================================================
\markboth{Índice de tablas}{Índice de tablas}
\listoftables
\thispagestyle{fancy}
\newpage

% ======================================================
% PÁGINA: ÍNDICE DE FIGURAS
% ======================================================
\markboth{Índice de figuras}{Índice de figuras}
\listoffigures
\thispagestyle{fancy}
\newpage

% ======================================================
% HOJA EN BLANCO CON MARCA DE AGUA
% ======================================================
\newpage
\pagestyle{fancy}

\begin{tikzpicture}[remember picture, overlay]
    \node[opacity=0.1] at (current page.center) {
\includegraphics[width=12cm]{logo_definitivo.png}};
\end{tikzpicture}

% Empujamos el texto hacia la parte inferior del folio
\vspace*{8.5cm} 
\begin{center}
    \textit{Esta hoja ha sido dejada deliberadamente en blanco.}
\end{center}
\vfill
\newpage

% ======================================================
% CUERPO DEL DOCUMENTO
% ======================================================

\section{DESARROLLO DEL PROYECTO}

\vspace{0.3cm}

% --- INTRODUCCIÓN A LA METODOLOGÍA ---
\begin{table}[H]
\centering\small
\begin{tabular}{|p{14.5cm}|}
\hline
\rowcolor{ArthroGray!15}
\textbf{\color{ArthroNavy}} Para que el desarrollo del proyecto se lleve a cabo de forma lineal y ordenada, se ha dividido el proceso en fases jerárquicas que permitan la retroalimentación y la minimización de errores en el diseño industrial, de forme que consigamos ahorrar costes y plazos temporales. A continuación, se definen las fases que componen la metodología de trabajo del Equipo 09. \\
\hline
\end{tabular}
\end{table}

\subsection{DIAGRAMA DE FASES}

\begin{table}[H]
\centering\small
\begin{tabular}{|p{14.5cm}|}
\hline
\end{tabular}
\end{table}

\begin{figure}[H]
    \centering
    \includegraphics[width=0.55\textwidth]{ANEXOS/diagrama_fases.png} 
    \caption{Diagrama de fases del desarrollo del producto.}
    \label{fig:diagrama_fases}
\end{figure}

\newpage
\subsection{DEFINICIÓN DE FASES}

\vspace{0.2cm}

% --- MATRIZ DE FASES ---
\begin{table}[H]
\centering\small
\begin{tabular}{|p{4.2cm}|p{9.8cm}|}
\hline
\rowcolor{ArthroNavy}
\textbf{\color{white}Fase de desarrollo} & \textbf{\color{white}Descripción y tareas asociadas} \\
\hline
\textbf{Plan estratégico y estudio de necesidades del usuario} & Se realiza una investigación cualitativa y cuantitativa mediante entrevistas personales y encuestas técnicas a pacientes con patologías respiratorias y articulares. El objetivo es recopilar datos sobre ergonomía de agarre y dificultades de presión para que el diseño se adhiera a los requisitos reales del usuario final. \\
\hline
\textbf{Definición de las especificaciones} & A partir de las necesidades detectadas, se establecen las Especificaciones de Diseño de Producto (EDP). Se desarrollan los 16 bloques de fuentes de especificaciones que actuarán como directrices fundamentales durante todo el proceso técnico y legal. \\
\hline
\textbf{Proceso de creación} & Esta fase constituye el núcleo creativo del desarrollo y se subdivide en tres etapas:
\begin{itemize}[leftmargin=*, topsep=0pt, itemsep=0.2em]
    \item \textbf{Diseño conceptual:} Generación de bocetos individuales por cada miembro del equipo para representar visualmente las soluciones de palanca y sujeción.
    \item \textbf{Diseño básico:} Evaluación de propuestas mediante matrices de decisión y selección del modelo base óptimo en términos de ergonomía y simplicidad.
    \item \textbf{Diseño en detalle:} Desarrollo del modelo final mediante herramientas de diseño asistido por ordenador (CAD). Se definen formalmente los tres componentes (cuerpo, palanca y base) y se generan los planos técnicos bajo norma.
\end{itemize} \\
\hline
\textbf{Validación} & Una vez finalizado el diseño detallado, se evalúa si el mecanismo de palanca cumple con el ratio de ventaja mecánica proyectado y si el ajuste sobre el inhalador es estable. En caso de discrepancia con las especificaciones iniciales, se procede a una revisión iterativa del diseño. \\
\hline
\textbf{Imagen de marca} & Se define la identidad visual de la empresa ArthroDesign Solutions S.L., de forma que se proyecte una imagen de dispositivos médicos accesibles, funcionales y de alta fiabilidad. \\
\hline
\textbf{Estudio de fabricación} & Se analizan los materiales adecuados (polímeros de grado médico) y los procesos de transformación (impresión 3D para prototipado e inyección para producción en serie), considerando la sostenibilidad y el ecodiseño. \\
\hline
\textbf{Estudio de presupuesto} & Se realiza un cálculo exhaustivo de costes, incluyendo materiales, mano de obra, gastos generales y beneficio industrial, de manera que podamos asegurar que el precio final por unidad sea competitivo en el sector de las ayudas técnicas. \\
\hline
\end{tabular}
\caption{Matriz de definición de fases de desarrollo del producto.}
\end{table}

\vspace{1.5cm}

\noindent\begin{minipage}{\textwidth}
    \begin{flushright}
        En Málaga, a \today
    \end{flushright}
    \vspace{2cm}

    \noindent
    \begin{minipage}[t]{0.45\textwidth}
        \centering
        \rule{6cm}{0.6pt} \\
        \textbf{Fdo. D. Sergio Carmona Muñoz} \\
        Colegiado Nº 06424 \\
        DNI: 26291930D
    \end{minipage}
    \hfill
    \begin{minipage}[t]{0.45\textwidth}
        \centering
        \rule{6cm}{0.6pt} \\
        \textbf{Fdo. D. Juan Gabriel Navarro Valdivielso} \\
        Colegiado Nº 06425 \\
        DNI: 77666662D
    \end{minipage}

    \vspace{2cm}

    \noindent
    \begin{minipage}[t]{0.45\textwidth}
        \centering
        \rule{6cm}{0.6pt} \\
        \textbf{Fdo. D. Antonio Sánchez Doña} \\
        Colegiado Nº 06423 \\
        DNI: 79380534J
    \end{minipage}
    \hfill
    \begin{minipage}[t]{0.45\textwidth}
        \centering
        \rule{6cm}{0.6pt} \\
        \textbf{Fdo. D. Iván de la Torre Segato} \\
        Colegiado Nº 06421 \\
        DNI: 79383630G
    \end{minipage}
\end{minipage}
\vspace{1cm}

% ======================================================
% HOJA EN BLANCO CON MARCA DE AGUA
% ======================================================
\newpage
\pagestyle{fancy}

\begin{tikzpicture}[remember picture, overlay]
    \node[opacity=0.1] at (current page.center) {
\includegraphics[width=12cm]{logo_definitivo.png}};
\end{tikzpicture}

% Empujamos el texto hacia la parte inferior del folio
\vspace*{8.5cm} 
\begin{center}
    \textit{Esta hoja ha sido dejada deliberadamente en blanco.}
\end{center}
\vfill

\end{document}